\chapter{Consentimiento}
\label{cap:consentimiento}

\chapter{Resumen}
\label{cap:resumen}
La interacción en videojuegos ha evolucionado más allá de los dispositivos convencionales como el
teclado y el ratón, especialmente en el ámbito de la rehabilitación y el ejercicio físico
(\textit{exergames}), que utilizan dispositivos de entrada especializados como sensores de
movimiento, cámaras de profundidad o mandos específicos para ofrecer una experiencia más natural
e intuitiva. Sin embargo, estos dispositivos especializados suelen ser costosos y específicos para
un tipo de juego o actividad, lo que limita su accesibilidad y uso general.

Este trabajo de fin de grado aborda dicha limitación mediante el diseño y desarrollo de un sistema de
entretenimiento y ejercicio que utiliza la detección de posturas del cuerpo humano a través de una cámara
como método de entrada principal.
Para ello, se utilizará un sistema de detección de posturas basado en el modelo de inteligencia artificial
MediaPipe BlazePose, que extrae la posición de puntos clave del cuerpo humano en tiempo real a partir de
imágenes capturadas por una cámara convencional. Este modelo es capaz de detectar hasta 33 puntos
articulares, así como generar una representación tridimensional de la postura del cuerpo humano.
Esta información será usada para controlar varios minijuegos desarrollados en el motor de videojuegos
Unity, que se ejecutarán en un navegador web a través de WebGL.

El resultado final es una plataforma web que permite al usuario mantener un registro de su forma y
actividad física, así como jugar a dos minijuegos distintos que utilizan la detección de posturas
para interactuar con el usuario. Esta aplicación web es multiplataforma, teniendo como único requisito
la utilización de un navegador web moderno y una cámara convencional. El sistema tiene la capacidad
de adaptar la calidad de la detección corporal a las capacidades de cómputo del dispositivo ejecutando
la aplicación. Gracias a esto, se consigue que el producto pueda ser utilizado en una amplia variedad
de dispositivos, desde ordenadores fijos hasta teléfonos móviles.

Los minijuegos desarrollados requieren que el usuario realice diferentes movimientos corporales,
como saltar, agacharse o intentar alcanzar un objeto con las manos, lo que proporciona una
experiencia de usuario lúdica y activa. En cada una de las sesiones de juego se registran tanto
la puntuación obtenida como la intensidad de actividad física realizada en forma de \textit{puntos
de actividad}. Esto permite al usuario llevar un control de la cantidad y calidad del ejercicio
realizado, todo desde la misma aplicación.

Este proyecto demuestra la viabilidad de crear un sistema de entretenimiento y ejercicio capaz de
detectar y analizar la postura del cuerpo humano en tiempo real utilizando únicamente una cámara
convencional, sin necesidad de dispositivos de entrada especializados ni de equipos de alto rendimiento.


\chapter{Abstract}
\label{cap:abstract}


\chapter{Agradecimientos}
\label{cap:agradecimientos}
